% --- Archon physics formalization source begin ---
% archon:physics
% archon:covers PhyXMiniProblems/problem_phyx_mini_0493.lean
% archon:source-report reports/phyx_mini/problem_phyx_mini_0493.source.json

\chapter{Physics problem phyx\_mini\_0493}
\label{ch:PhyXMiniProblems_problem_phyx_mini_0493}

\paragraph{Problem source.}
\#\# Physical scenario

A heat engine's high temperature \textbackslash{}(T\_H\textbackslash{}) could be ambient temperature, because liquid nitrogen at \textbackslash{}(77\textbackslash{},\textbackslash{}mathrm\{K\}\textbackslash{}) could be \textbackslash{}(T\_L\textbackslash{}) and is cheap.The Carnot engine made use of heat transferred from air at room temperature (\textbackslash{}(293\textbackslash{},\textbackslash{}mathrm\{K\}\textbackslash{})) to the liquid nitrogen fuel (figure).

\#\# Question

What would be the efficiency of a Carnot engine?

\#\# Answer choices

- A: A: 68\%

- B: B: 70\%

- C: C: 72\%

- D: D: 74\%

\#\# Figure caption (auxiliary; use the image as primary evidence)

The image depicts a laboratory or industrial setting with a person in a lab coat standing near large metal cryogenic tanks. There are four tanks visible, labeled with numbers "0," "2," "3," and the letter "K." The tanks are emitting vapor, indicating they contain a cold substance, likely liquid nitrogen. The person is looking at a sheet of paper and is wearing gloves, suggesting safety precautions. Next to the person is a metal cart with a clipboard and paperwork. The environment appears clean and organized, suggesting a controlled scientific or industrial process.

\#\# Dataset metadata

Laws of Thermodynamics; Physical Model Grounding Reasoning, Multi-Formula Reasoning

\paragraph{Recorded answer/context.}
D: D: 74\%

\paragraph{Figure/image path.}
/root/proposal\_for\_physic/science-mango/phyx\_mini\_run/phyx\_data/test\_image/493.png

\paragraph{Formalization target.}
create a compiling Lean file with sorry bodies at `PhyXMiniProblems/problem\_phyx\_mini\_0493.lean`. The Lean declarations must preserve the physical quantities, dimensions or dimensional roles, figure labels, governing-law hypotheses, and final relation expressed by this problem.
Use Mathlib/Physlib names found through LeanExplore where available. If a physics API is missing, introduce faithful local abstractions rather than scalar placeholder aliases.

\begin{theorem}[Physics formalization target]
\label{thm:physics:phyx_mini_0493:target}
\lean{PhyXMiniProblems.ProblemPhyXMini0493.carnotEfficiency_rounds_to_recordedAnswerD}
\uses{def:physics:phyx-mini-0493:phyxminiproblems-problemphyxmini0493-ambientnitrogencarnotsetup, def:physics:phyx-mini-0493:phyxminiproblems-problemphyxmini0493-matchesproblemstatement, def:physics:phyx-mini-0493:phyxminiproblems-problemphyxmini0493-matchesprimarycryogenicfigure, def:physics:phyx-mini-0493:phyxminiproblems-problemphyxmini0493-hasphysicalcarnotparameters, def:physics:phyx-mini-0493:phyxminiproblems-problemphyxmini0493-satisfiescarnotefficiencylaw, def:physics:phyx-mini-0493:phyxminiproblems-problemphyxmini0493-recordedanswerchoice}
The assigned autoformalize agent should translate this physics problem into Lean declarations in the covered file, with theorem and lemma proof bodies written as `by sorry`.
\end{theorem}
\begin{proof}
This is an autoformalization task, not a proof task. Produce statements that can later be proved by the physics prover without weakening the physical model.
\end{proof}

% --- Archon declaration topology begin ---
\section{Declaration topology}
\begin{definition}[Absolute Temperature Quantity]
  \label{def:physics:phyx-mini-0493:phyxminiproblems-problemphyxmini0493-absolutetemperaturequantity}
  \lean{PhyXMiniProblems.ProblemPhyXMini0493.AbsoluteTemperatureQuantity}
  A nonnegative absolute thermodynamic temperature
\end{definition}
\begin{proof}
  This is the declared model definition of Absolute Temperature Quantity.
\end{proof}
\begin{definition}[temperature In Kelvins]
  \label{def:physics:phyx-mini-0493:phyxminiproblems-problemphyxmini0493-temperatureinkelvins}
  \lean{PhyXMiniProblems.ProblemPhyXMini0493.temperatureInKelvins}
  \uses{def:physics:phyx-mini-0493:phyxminiproblems-problemphyxmini0493-absolutetemperaturequantity}
  Read an absolute temperature in kelvins from its stated storage unit
\end{definition}
\begin{proof}
  This is the declared model definition of temperature In Kelvins.
\end{proof}
\begin{definition}[Thermodynamic Device Role]
  \label{def:physics:phyx-mini-0493:phyxminiproblems-problemphyxmini0493-thermodynamicdevicerole}
  \lean{PhyXMiniProblems.ProblemPhyXMini0493.ThermodynamicDeviceRole}
  The thermodynamic role of the device described in the problem
\end{definition}
\begin{proof}
  This is the declared model definition of Thermodynamic Device Role.
\end{proof}
\begin{definition}[Heat Engine Model]
  \label{def:physics:phyx-mini-0493:phyxminiproblems-problemphyxmini0493-heatenginemodel}
  \lean{PhyXMiniProblems.ProblemPhyXMini0493.HeatEngineModel}
  The reversible ideal-engine model named by the question
\end{definition}
\begin{proof}
  This is the declared model definition of Heat Engine Model.
\end{proof}
\begin{definition}[Reservoir]
  \label{def:physics:phyx-mini-0493:phyxminiproblems-problemphyxmini0493-reservoir}
  \lean{PhyXMiniProblems.ProblemPhyXMini0493.Reservoir}
  The two thermal reservoirs used in the Carnot model
\end{definition}
\begin{proof}
  This is the declared model definition of Reservoir.
\end{proof}
\begin{definition}[Reservoir Substance]
  \label{def:physics:phyx-mini-0493:phyxminiproblems-problemphyxmini0493-reservoirsubstance}
  \lean{PhyXMiniProblems.ProblemPhyXMini0493.ReservoirSubstance}
  The material supplying each reservoir in the narrated scenario
\end{definition}
\begin{proof}
  This is the declared model definition of Reservoir Substance.
\end{proof}
\begin{definition}[Cryogenic Vessel Mark]
  \label{def:physics:phyx-mini-0493:phyxminiproblems-problemphyxmini0493-cryogenicvesselmark}
  \lean{PhyXMiniProblems.ProblemPhyXMini0493.CryogenicVesselMark}
  Literal marks visibly painted on the cryogenic vessels in the bitmap
\end{definition}
\begin{proof}
  This is the declared model definition of Cryogenic Vessel Mark.
\end{proof}
\begin{definition}[Cryogenic Facility Figure]
  \label{def:physics:phyx-mini-0493:phyxminiproblems-problemphyxmini0493-cryogenicfacilityfigure}
  \lean{PhyXMiniProblems.ProblemPhyXMini0493.CryogenicFacilityFigure}
  \uses{def:physics:phyx-mini-0493:phyxminiproblems-problemphyxmini0493-cryogenicvesselmark}
  Qualitative and literal data visible in the primary bitmap. These fields keep the supplied image in the formal model even though the efficiency calculation uses the temperatures stated in the prose
\end{definition}
\begin{proof}
  This is the declared model definition of Cryogenic Facility Figure.
\end{proof}
\begin{definition}[Ambient Nitrogen Carnot Setup]
  \label{def:physics:phyx-mini-0493:phyxminiproblems-problemphyxmini0493-ambientnitrogencarnotsetup}
  \lean{PhyXMiniProblems.ProblemPhyXMini0493.AmbientNitrogenCarnotSetup}
  \uses{def:physics:phyx-mini-0493:phyxminiproblems-problemphyxmini0493-absolutetemperaturequantity, def:physics:phyx-mini-0493:phyxminiproblems-problemphyxmini0493-thermodynamicdevicerole, def:physics:phyx-mini-0493:phyxminiproblems-problemphyxmini0493-heatenginemodel, def:physics:phyx-mini-0493:phyxminiproblems-problemphyxmini0493-reservoir, def:physics:phyx-mini-0493:phyxminiproblems-problemphyxmini0493-reservoirsubstance, def:physics:phyx-mini-0493:phyxminiproblems-problemphyxmini0493-cryogenicfacilityfigure}
  The physical setup. The efficiency is an unconstrained dimensionless field until the generic Carnot governing law is imposed below
\end{definition}
\begin{proof}
  This is the declared model definition of Ambient Nitrogen Carnot Setup.
\end{proof}
\begin{definition}[Matches Problem Statement]
  \label{def:physics:phyx-mini-0493:phyxminiproblems-problemphyxmini0493-matchesproblemstatement}
  \lean{PhyXMiniProblems.ProblemPhyXMini0493.MatchesProblemStatement}
  \uses{def:physics:phyx-mini-0493:phyxminiproblems-problemphyxmini0493-temperatureinkelvins, def:physics:phyx-mini-0493:phyxminiproblems-problemphyxmini0493-ambientnitrogencarnotsetup}
  Problem-statement data: heat is taken from ambient air at room temperature and transferred to an inexpensive liquid-nitrogen sink. These premises do not assign a value to the requested efficiency
\end{definition}
\begin{proof}
  This is the declared model definition of Matches Problem Statement.
\end{proof}
\begin{definition}[Matches Primary Cryogenic Figure]
  \label{def:physics:phyx-mini-0493:phyxminiproblems-problemphyxmini0493-matchesprimarycryogenicfigure}
  \lean{PhyXMiniProblems.ProblemPhyXMini0493.MatchesPrimaryCryogenicFigure}
  \uses{def:physics:phyx-mini-0493:phyxminiproblems-problemphyxmini0493-cryogenicvesselmark, def:physics:phyx-mini-0493:phyxminiproblems-problemphyxmini0493-ambientnitrogencarnotsetup}
  Primary-image evidence. The raster visibly contains marks 1, 2, 3, and K, metal cryogenic vessels emitting vapor, a lab-coated attendant, and a paperwork cart. No efficiency value is read from the image
\end{definition}
\begin{proof}
  This is the declared model definition of Matches Primary Cryogenic Figure.
\end{proof}
\begin{definition}[Has Physical Carnot Parameters]
  \label{def:physics:phyx-mini-0493:phyxminiproblems-problemphyxmini0493-hasphysicalcarnotparameters}
  \lean{PhyXMiniProblems.ProblemPhyXMini0493.HasPhysicalCarnotParameters}
  \uses{def:physics:phyx-mini-0493:phyxminiproblems-problemphyxmini0493-temperatureinkelvins, def:physics:phyx-mini-0493:phyxminiproblems-problemphyxmini0493-ambientnitrogencarnotsetup}
  Positivity, ordering, and range conditions for a physical heat engine
\end{definition}
\begin{proof}
  This is the declared model definition of Has Physical Carnot Parameters.
\end{proof}
\begin{definition}[Satisfies Carnot Efficiency Law]
  \label{def:physics:phyx-mini-0493:phyxminiproblems-problemphyxmini0493-satisfiescarnotefficiencylaw}
  \lean{PhyXMiniProblems.ProblemPhyXMini0493.SatisfiesCarnotEfficiencyLaw}
  \uses{def:physics:phyx-mini-0493:phyxminiproblems-problemphyxmini0493-temperatureinkelvins, def:physics:phyx-mini-0493:phyxminiproblems-problemphyxmini0493-ambientnitrogencarnotsetup}
  The governing law for a reversible Carnot engine. It is stated generically in terms of the two reservoir temperatures and does not contain the requested numerical efficiency or answer choice
\end{definition}
\begin{proof}
  This is the declared model definition of Satisfies Carnot Efficiency Law.
\end{proof}
\begin{definition}[Answer Choice]
  \label{def:physics:phyx-mini-0493:phyxminiproblems-problemphyxmini0493-answerchoice}
  \lean{PhyXMiniProblems.ProblemPhyXMini0493.AnswerChoice}
  Labels attached to the four displayed multiple-choice answers
\end{definition}
\begin{proof}
  This is the declared model definition of Answer Choice.
\end{proof}
\begin{definition}[efficiency Percent]
  \label{def:physics:phyx-mini-0493:phyxminiproblems-problemphyxmini0493-answerchoice-efficiencypercent}
  \lean{PhyXMiniProblems.ProblemPhyXMini0493.AnswerChoice.efficiencyPercent}
  \uses{def:physics:phyx-mini-0493:phyxminiproblems-problemphyxmini0493-answerchoice}
  Whole efficiency percentage printed beside each answer label
\end{definition}
\begin{proof}
  This is the declared model definition of efficiency Percent.
\end{proof}
\begin{definition}[recorded Answer Choice]
  \label{def:physics:phyx-mini-0493:phyxminiproblems-problemphyxmini0493-recordedanswerchoice}
  \lean{PhyXMiniProblems.ProblemPhyXMini0493.recordedAnswerChoice}
  \uses{def:physics:phyx-mini-0493:phyxminiproblems-problemphyxmini0493-answerchoice}
  The answer label recorded in the dataset metadata
\end{definition}
\begin{proof}
  This is the declared model definition of recorded Answer Choice.
\end{proof}
\begin{lemma}[carnot Efficiency eq 216 div 293]
  \label{lem:physics:phyx-mini-0493:phyxminiproblems-problemphyxmini0493-carnotefficiency-eq-216-div-293}
  \lean{PhyXMiniProblems.ProblemPhyXMini0493.carnotEfficiency_eq_216_div_293}
  \uses{def:physics:phyx-mini-0493:phyxminiproblems-problemphyxmini0493-ambientnitrogencarnotsetup, def:physics:phyx-mini-0493:phyxminiproblems-problemphyxmini0493-matchesproblemstatement, def:physics:phyx-mini-0493:phyxminiproblems-problemphyxmini0493-matchesprimarycryogenicfigure, def:physics:phyx-mini-0493:phyxminiproblems-problemphyxmini0493-hasphysicalcarnotparameters, def:physics:phyx-mini-0493:phyxminiproblems-problemphyxmini0493-satisfiescarnotefficiencylaw}
  The exact Carnot efficiency is 1 - 77 / 293 = 216 / 293
\end{lemma}
\begin{proof}
  The conclusion follows from the governing relations and figure readouts named in the statement of carnot Efficiency eq 216 div 293.
\end{proof}
% --- Archon declaration topology end ---
% --- Archon physics formalization source end ---
